\section{Introduction}

Multi-agent AI systems that recursively delegate tasks face a structural failure mode we call the \textbf{autonomous drift problem}. As an agent spawns sub-agents that spawn further sub-agents, delegation chains lose traceable connectivity to the original authorizing principal. Orphaned agents continue operating indefinitely; drifted agents execute objectives that no human ever endorsed; and in the worst case, both conditions propagate simultaneously down the spawn chain through what we term a \textbf{drift cascade} — a single undetected drift event at depth $d$ with branching factor $b$ producing up to $b^d$ drifted agents in the field. This is not a governance concern reducible to policy conventions. It is a correctness failure: a system that has departed from its original intent with no structural mechanism to detect or recover that departure.

The field has attempted to address drift through two primary mechanisms. \textbf{Time-to-live} (TTL) constraints limit the operational lifetime of an agent or the number of delegation hops before workflow termination. \textbf{Depth limits} cap the maximum nesting depth of recursive spawning. Both approaches impose an external, temporal bound on system behavior. Neither addresses the fundamental structural problem: when a sub-agent spawns a further sub-agent, the system retains no immutable, queryable record of the parent-child relationship. TTL expiration does not preserve lineage — it terminates. Depth limits do not encode intent propagation — they truncate. As a result, existing systems exhibit a structural blind spot: they can enforce \emph{when} a system stops, but not \emph{what} it did or \emph{why}. The deeper structural reason depth limits are insufficient is that they address the wrong dimension of the accountability problem. Depth limits constrain chain length — a topological property of the spawn graph. Drift, however, is fundamentally about authorization scope: whether each node's operating envelope is a legitimate descendant refinement of the human anchor's intent. A drifted agent at depth 1 with unbounded scope expansion authority can produce consequences that exceed the human anchor's intent as severely as a drifted agent at depth $D_{\max} - 1$. We prove formally in Section~\ref{sec:autonomous-drift} that no choice of $D_{\max}$ prevents scope creep drift without an additional scope refinement constraint.

The core insight of this paper is that the autonomous drift problem is structurally soluble: it is not an inherent consequence of recursive spawning, but a consequence of mutable parent pointers. If every agent's parent reference is fixed immutably at spawn time and embedded in the agent's foundational identity — not stored as a mutable attribute — then the agent cannot dissociate from its chain of authorization. We formalize this through the \textbf{immutable parent pointer chain}: a cryptographically signed, append-only chain of ancestor references sealed at agent creation via SHA-256 forensic sealing. Modification of any ancestor record invalidates the chain's integrity. This makes the drift triad — the conjunction of orphaning, drift from human anchor, and continued operation — architecturally impossible, not merely statistically unlikely. An agent cannot be orphaned from a parent it structurally cannot forget. It cannot be drifted from an intent that is cryptographically embedded in its identity. And it cannot operate without having passed all three BX3 layer validations at activation time.

BX3 operationalizes immutable parent chains through the \textbf{Behavior, Execution, and eXchange} (BEX) architecture. Every agent maintains a \textbf{Parent Pointer Chain} (PPC) — a provably complete, tamper-evident ancestry record sealed at provisioning. The \textbf{Chain Query Interface} (CQI) enables any authorized agent or human overseer to reconstruct the full delegation path at any depth. The \textbf{Purpose Layer} serves as the singular, immutable anchor: every root agent is initialized with a declared \textit{purpose clause} that defines the boundaries of its operational scope. No child agent may operate outside the semantic scope defined by its parent's purpose, and no agent may revise its purpose post-spawn. The \textbf{Bounds Engine} enforces structural and computational constraints — maximum depth $D_{\max}$ and resource budgets — preventing unbounded recursion. The \textbf{Fact Layer} is the append-only, tamper-evident ledger that records every spawn event, delegation action, purpose clause binding, and agent state transition, providing the forensic guarantee central to all correctness claims: the complete ancestry of any agent is recoverable with cryptographic certainty.

\textbf{Formal definition of drift.} We define drift formally (Section~\ref{sec:autonomous-drift}) as the conjunction of three independent conditions constituting the \textbf{drift triad}:
\[
\mathrm{DriftTriad}(n) = \langle \mathrm{orphaned}(n),\; \mathrm{drifted}(n),\; \mathrm{operating}(n) \rangle.
\]
An agent is \emph{orphaned} when its parent has terminated or become unreachable. An agent is \emph{drifted} when its delegation chain does not reach any human principal at its origin. An agent \emph{operates} when it continues executing its event loop. The drift triad — all three conditions holding simultaneously — is the specific failure state the Recursive Spawning Protocol targets. Drift is heritable: a drifted agent that spawns a child propagates the condition to all descendants (Equation~\ref{eq:drift-cascade}), producing exponential amplification. The drift prevention requirement is precisely the conjunction of two independent constraints:
\[
\mathrm{Drift\ Prevention} \iff \bigl( \mathrm{depth}(C) \leq D_{\max} \bigr)\ \land\ \bigl( \forall i:\ R(n_{i+1}) \subseteq R(n_i) \bigr),
\]
where the first conjunct is the depth constraint and the second is the scope refinement constraint. Depth limits alone enforce only the first. BX3 enforces both structurally at the Fact Layer pre-commit hook.

\textbf{Paper roadmap.} The remainder of this paper is organized as follows. Section~\ref{sec:autonomous-drift} formalizes the drift problem, establishes the drift triad and its constituent sub-modes, identifies three structural enabling conditions, catalogs four concrete failure modes, and proves that depth limits alone are insufficient. Section~\ref{sec:bx3-integration} presents the integration of Recursive Spawning with the BX3 Purpose Layer, Bounds Engine, and Fact Layer, showing how they collectively enforce drift prevention as a structural invariant. Section~\ref{sec:rs-protocol} specifies the Recursive Spawning Protocol (RSP) — the mandatory five-stage spawn sequence executed across all three BX3 layers at every spawn event — and proves three correctness properties: drift prevention, chain integrity, and termination. Section~\ref{sec:experimental} presents the experimental evaluation measuring drift resistance under adversarial propagation scenarios. Sections~\ref{sec:limitations} and~\ref{sec:conclusion} discuss limitations, ethical implications, and open problems.